\documentclass[12pt]{article}
\usepackage[utf8]{inputenc}
\usepackage{amsmath, amssymb,multicol}
\usepackage{hyperref}
\usepackage{tikz}
\usepackage[left=2.5cm,right=2.5cm,bottom=2.5cm]{geometry}
\usetikzlibrary{arrows, arrows.meta, positioning, bending, quotes}
\hypersetup{
    colorlinks=true}

\title{MATH 5707 (Spring 2026): Homework 9}
\date{}

\begin{document}
\vspace{-0.7in}

\maketitle
\vspace{-0.6in}


\section*{Directions and Introduction}
Submit a pdf of your solutions to the HW 9 assignment on Gradescope by 11:59 pm on April 10.

Problem 3 is marked as a peer review problem. \href{https://docs.google.com/document/d/1jSw9pmMJJFUx_6dTkcUi4k4HJNIrrIIgs9zdWhAycx0/edit?usp=sharing}{This document gives directions and deadlines for the peer review process.}



\section*{Problems}

\begin{enumerate}
\item[0.] If you would like any of these problems to be graded for proficiency with the key skills, list the skill and the corresponding problem.  If you want to label the skill by number, use the numbering that is in the syllabus.

\item Draw a graph $G$ that satisfies the following properties, or explain why no such graph can exist.
\begin{enumerate}
\item a graph $G$ with $\Delta(G) = 2\chi (G)$
\item a graph $G$ with $\chi (G) = 2 \Delta (G)$
\item a noncomplete graph of order $n$ with chromatic number $n$
\end{enumerate}

    \item(Each part is worth 4 points and
graded as a separate problem.) Prove or disprove.

\begin{enumerate}
\item If $G$ is a simple graph with $\chi(G)\ge3$, then $K_3\subseteq G$. 
\item If $G$ is an $r$-regular graph with an odd number of vertices, then $\chi'(G)=r+1$.
\end{enumerate}
   
   \item (Peer review problem; Bondy-Murty 8.1.10(b)) A graph $G$ is called \emph{critical} if $\chi(H)<\chi(G)$ for all proper subgraphs $H\subset G$. Prove that $G_1+G_2$ is critical if and only if both $G_1$ and $G_2$ are critical. (Note: You can use the fact that we proved part (a) of the problem as an activity in class.)
   
   
       \item The Math Department at Complex U. has eight faculty members, each of whom is teaching three courses this semester. The course numbers are 127, 131, 132, 136, 138, 153, 154, 201, 205, 211. The schedules are:

\begin{multicols}{2} 
Agnesi 132, 136, 211 

Bernoulli 127, 131, 153 

Cauchy 131, 132, 211 

Descartes 127, 131, 205

Euler 131, 138, 154

Frobenius 132, 136, 201

Gauss 127, 131, 138

Hamilton 153, 154, 205
\end{multicols}

It has been decided that each course will have a department-wide final exam such that all sections of a given course will have their finals at the same time. Each professor must proctor their own exam. There is plenty of classroom space available, but everyone would like to get the exams over as soon as possible so they can rush off to winter break and prove lots of new theorems. What is the fewest number of time slots that can be used to give the exams?

\begin{enumerate}
\item To set up this problem as a vertex coloring problem, we can let...

\begin{itemize}
\item Vertices represent \rule{1cm}{0.15mm}
\item Adjacent vertices represent \rule{1cm}{0.15mm}
\item Colors of vertices represent \rule{1cm}{0.15mm}
\item The chromatic number represents \rule{1cm}{0.15mm}
\end{itemize}
\item Draw an appropriate graph and use vertex coloring to find an exam schedule that requires the least number of time slots. Be sure to justify how you know you have found the minimum. 
\end{enumerate}
    
  
\end{enumerate}

\end{document}